\begin{abstract}
We study repeated queries over stable documents. \textbf{CoMem} stores one
depth-$j$ residual per token, retrieves a bounded chunk set, and resumes layers
$[j{:}L)$. In a matched comparison with identical evidence, mask, examples, and
adapter, $j{=}12$ reduces isolated model Read from 931.9 to 664.4\,ms
($1.403\times$) while lowering a 15-cell RULER macro from 99.19 to 96.07
(paired gap 3.12, 95\% CI $[2.36,3.93]$). A rank-32 self-distillation adapter
makes the interface usable without updating the backbone. The 8\,KiB/token
store requires reuse: at 32k, CPU-pinned break-even is 8.9--10.9 queries for
generations up to 128 tokens; at 128k and one generated token it is
25.8--27.6 across GPU/CPU placement. In a nine-cell equal-latency diagnostic,
dependence-aware hierarchical intervals show BM25 replay ahead by 11.56 points
($[-18.67,-5.11]$), while frozen-BGE replay has an unresolved 1.00-point edge
($[-10.67,8.33]$); CoMem wins neither aggregate. Restoring lower-layer document
context raises a paired multikey diagnostic from 92.5 to 100.0, and a 32-token
overlap reaches 98.5 without increasing persistent bytes or per-query Read.
CoMem is an internal measurement of a depth-reuse operating point, not a claim
of superiority over raw-text retrieval or modular-cache systems.
\end{abstract}
